% Teil II: der durchgerechnete Weg.

% Gedankenschritt: Kürzen durch die führende Potenz/Basis ist der Schlüssel;
% die kleineren Terme sterben für n->infty ab.

\subsection*{Gegeben und gesucht}
Grenzwerte der Folgen $a_n,b_n,c_n,d_n$ f\"ur $n\to\infty$.

\subsection*{Ansatz}
F\"ur jeden Quotienten: die \emph{gemeinsame f\"uhrende Ordnung} im
Z\"ahler und Nenner finden und durch sie k\"urzen. Damit bleiben nur
Terme, deren Grenzwert man direkt (zu $0$ oder zu einer Konstante) sieht.

\subsection*{a) $a_n=\frac{3n^3-2n^2+1}{n^2-2n^3-3}$}
Die h\"ochste Potenz ist $n^3$, im Z\"ahler wie Nenner. K\"urzen durch
$n^3$:
\begin{equation}
  a_n=\frac{3-\frac2n+\frac1{n^3}}{\frac1n-2-\frac3{n^3}}
  \;\to\; \frac{3-0+0}{0-2-0}=-\frac32 .
\end{equation}
Gedanke: $\tfrac1n,\tfrac1{n^3}\to0$; \"uberleben nur $3$ und $-2$.

\subsection*{b) $b_n=\frac{\sqrt n+2}{\sqrt{n+1}-1}$}
Die f\"uhrende Ordnung ist $\sqrt n$. Durch $\sqrt n$ k\"urzen
($\tfrac{1}{\sqrt n},\tfrac1n\to0$):
\begin{equation}
  b_n=\frac{1+\frac2{\sqrt n}}{\sqrt{1+\frac1n}-\frac1{\sqrt n}}
  \;\to\;\frac{1+0}{1-0}=1 .
\end{equation}
Gedanke: $\sqrt{n+1}=\sqrt{n}\sqrt{1+\tfrac1n}\approx\sqrt n$ f\"ur
gro\ss{}e $n$.

\subsection*{c) $c_n=\frac{-2\cdot3^n+(-1)^n}{5\cdot2^n-4\cdot3^{n+1}}$}
$3^n$ w\"achst schneller als $2^n$, und $3^{n+1}=3\cdot3^n$; die gr\"o\ss{}te
Basis ist also $3^n$. K\"urzen durch $3^n$:
\begin{equation}
  c_n=\frac{-2+\frac{(-1)^n}{3^n}}{5\cdot\bigl(\frac23\bigr)^n-12}
  \;\to\;\frac{-2+0}{0-12}=\frac16 .
\end{equation}
Gedanken: $\frac{(-1)^n}{3^n}\to0$ (beschr\"ankter Z\"ahler, explodierender
Nenner), $(\tfrac23)^n\to0$, $-4\cdot3^{n+1}/3^n=-12$.

\subsection*{d) $d_n=\frac{n^2}{(n+1)!}$}
Es ist $(n+1)!=(n+1)\,n!$; die Fakult\"at w\"achst schneller als jede
feste Potenz:
\begin{equation}
  d_n=\frac{n^2}{(n+1)\,(n!)}\;\to\;0 .
\end{equation}
Gedanke: Schreibt man $(n+1)!$ aus, steht $n!$ im Nenner und der Z\"ahler
ist eine fixe Potenz --- der Quotient schrumpft gegen $0$.

\begin{ergebnis}
\begin{equation*}
  \text{a)}\ -\frac32,\qquad
  \text{b)}\ 1,\qquad
  \text{c)}\ \frac16,\qquad
  \text{d)}\ 0 .
\end{equation*}
\end{ergebnis}

\begin{kontrolle}
\textbf{Numerisch (f\"ur gro\ss{}e $n$):} $n=10^6$: $a_n\approx-1.5$,
$b_n\approx1$, $c_n\approx\frac16$, $d_n$ winzig. Die f\"uhrenden
Ordnungen best\"atigen die Werte:
$a_n\sim\frac{3n^3}{-2n^3}$, $b_n\sim\frac{\sqrt n}{\sqrt n}$, $c_n\sim
\frac{-2\cdot3^n}{-12\cdot3^n}$, $d_n\sim\frac{n^2}{n!}$.

\textbf{Grenzfall:} W\"urde man in a) durch $n^2$ statt $n^3$ k\"urzen,
w\"urde der Nenner $\frac1n-2-\dots$ zwar nach $-2$ gehen, aber der
Z\"ahler $\to\infty$ --- man m\"usste erst die h\"ochste Potenz statt
einer beliebigen w\"ahlen. Genau das ist der Kern der f\"uhrenden
Ordnung: Das Verfahren liefert sinnvolle Werte nur mit der gr\"o\ss{}ten
Ordnung.
\end{kontrolle}

\begin{merksatz}
Mit der f\"uhrenden Potenz bzw. der gr\"o\ss{}ten Basis k\"urzen;
Potenzen-, Wurzel- und Fakult\"atsangaben dabei sauber umschreiben
($\sqrt{n+1}\approx\sqrt n$, $3^{n+1}=3\cdot3^n$). Die kleineren Terme
sterben f\"ur $n\to\infty$ ab.
\end{merksatz}
